\paragraph{常数 \texorpdfstring{$4$}{4} 的刚性刻画：残数标尺、深度权重与饱和标尺的三重一致化}
在前述视界接口中，“\(4\)”并非引入新的动力学常数，而是由制度闭包的三条独立一致性要求所共同锁定：
其一是几何侧的曲率标尺（命题 \ref{prop:dephys-hyperbolic-curvature-normalization-four}）；
其二是解析侧的残数标尺（下述命题 \ref{prop:xi-four-rigidity-residue-depth-normalization}）；
其三是计数侧的饱和标尺（命题 \ref{prop:xi-fourpi-entropy-normalization-uniqueness}）。
本段把这三条要求在扫描层析的标量证书语言中严格化。

\begin{proposition}[系数 \texorpdfstring{$4$}{4} 的刚性刻画：谱残数与几何深度的一致化唯一确定规一化常数]\label{prop:xi-four-rigidity-residue-depth-normalization}
\leanverified{paper\_xi\_four\_rigidity\_residue\_depth\_normalization}
在有限点状缺陷口径下，令共动剖面密度取一族待定规一化
\[
H_c(x)
:=\sum_{j=1}^{\kappa} m_j\,\frac{c\,\delta_j}{(x-\gamma_j)^2+(1+\delta_j)^2},
\qquad c>0.
\]
令其 Cauchy 变换（\(\Im z>0\)）
\[
M_{H_c}(z):=\int_{\RR}\frac{H_c(x)}{x-z}\,\dd x,
\]
并保持内函数 \(S_{\mathrm{HP}}\) 与
\(
M_P(z):=-2\mathrm{i}\,\frac{d}{dz}\log S_{\mathrm{HP}}(z)
\)
的约定（见定理 \ref{thm:xi-double-herglotz-residue-ratio}）。
则二者在下半平面极点 \(z_j^-=\gamma_j-\mathrm{i}(1+\delta_j)\) 的残数满足
\[
\operatorname*{Res}_{z=z_j^-} M_{H_c}(z)=-\pi c\,m_j\,\frac{\delta_j}{1+\delta_j},
\qquad
\operatorname*{Res}_{z=z_j^-} M_{P}(z)=2\mathrm{i}\,m_j,
\]
从而
\[
\frac{1}{2\pi}\left|\frac{\operatorname*{Res}_{z=z_j^-} M_{H_c}(z)}{\operatorname*{Res}_{z=z_j^-} M_P(z)}\right|
=\frac{c}{4}\,\frac{\delta_j}{1+\delta_j}.
\]
若要求“由残数比恢复的无量纲深度权重”与几何 Busemann 深度权重 \(\delta_j/(1+\delta_j)\) 完全一致，即
\[
\frac{\delta_j}{1+\delta_j}
=\frac{1}{2\pi}\left|\frac{\operatorname*{Res}_{z=z_j^-} M_{H_c}(z)}{\operatorname*{Res}_{z=z_j^-} M_P(z)}\right|
\quad(\forall j),
\]
则必有且仅有 \(c=4\)。
\end{proposition}

\begin{proof}
对单原子 \(H_{c,(\gamma,\delta,m)}(x):=m\,\frac{c\delta}{(x-\gamma)^2+w^2}\)（\(w=1+\delta\)）重复定理 \ref{thm:xi-double-herglotz-residue-ratio} 证明中的留数计算（只需把 \(4\) 替换为 \(c\)）得到
\[
M_{H_c}(z)
=-\pi c\,m\,\frac{\delta}{w}\cdot\frac{1}{z-(\gamma-\mathrm{i}w)},
\]
故在 \(z^-=\gamma-\mathrm{i}(1+\delta)\) 处残数为 \(-\pi c\,m\,\frac{\delta}{1+\delta}\)，多原子情形按线性叠加即得第一式。
第二式由定理 \ref{thm:xi-double-herglotz-residue-ratio} 已证。
将二者代入并取绝对值即可得到 \(\frac{c}{4}\frac{\delta}{1+\delta}\)；因此要求对一切 \(j\) 成立只能强制 \(c=4\)。证毕。
\end{proof}

\begin{proposition}[深度权重的唯一性：乘法伸缩半群的 Möbius 线性化强制 \texorpdfstring{$p(\delta)=\delta/(1+\delta)$}{p(delta)=delta/(1+delta)}]\label{prop:xi-depth-weight-uniqueness-mobius-linearization}
\leanverified{paper\_xi\_depth\_weight\_uniqueness\_mobius\_linearization}
考虑深度伸缩作用 \(\delta\mapsto m\delta\)（\(m>0\)）。
设存在连续严格单调函数 \(p:(0,\infty)\to(0,1)\)，使得对所有 \(m>0\) 与 \(\delta>0\) 有
\[
p(m\delta)=\Phi_m\!\bigl(p(\delta)\bigr),
\qquad
\Phi_m(p):=\frac{m p}{1+(m-1)p},
\]
其中 \(\Phi_m\) 为保持端点 \(0,1\) 的分式线性变换。
则必存在常数 \(a>0\) 使得
\[
\boxed{\;
p(\delta)=\frac{\delta}{\delta+a}\;}
\]
并且在本文 unitary-slice 规一化下（把 \(a\) 固定为 \(1\)）得到唯一形式
\[
\boxed{\;
p(\delta)=\frac{\delta}{1+\delta}\; }.
\]
\end{proposition}

\begin{proof}
令 \(q(\delta):=\frac{p(\delta)}{1-p(\delta)}\in(0,\infty)\)。
由函数方程直接计算得到
\[
q(m\delta)=\frac{p(m\delta)}{1-p(m\delta)}=\frac{\Phi_m(p(\delta))}{1-\Phi_m(p(\delta))}=m\,q(\delta).
\]
由连续性可知 \(q(\delta)=\delta/a\)（\(a:=q(1)^{-1}>0\)），从而
\(
p(\delta)=\frac{q(\delta)}{1+q(\delta)}=\frac{\delta}{\delta+a}
\)。
本文约定的 Cayley/unitary-slice 标尺把“单位深度”选为 \(\delta=1\) 时的对称点，从而 \(p(1)=\tfrac12\) 并固定 \(a=1\)。
证毕。
\end{proof}

\begin{proposition}[边界作用量表述：缺陷熵是 Green 势的加权 Neumann 边界项，规一化常数被唯一锁定]\label{prop:xi-defect-entropy-green-neumann-boundary-action}
令 \(t_j:=\gamma_j+\mathrm{i}(1+\delta_j)\in\mathbb H\)，并定义 Green 型势
\[
U_j(t):=\log\left|\frac{t-\overline{t_j}}{t-t_j}\right|
\qquad(t\in\mathbb H).
\]
则其在边界 \(y=0\) 的法向导数满足
\[
\partial_y U_j(x)=\frac{2(1+\delta_j)}{(x-\gamma_j)^2+(1+\delta_j)^2}.
\]
进一步定义加权边界作用量
\[
\mathcal S
:=\frac{1}{4\pi}\int_{\RR}\sum_{j=1}^{\kappa}m_j\,\frac{2\delta_j}{1+\delta_j}\,\partial_y U_j(x)\,\dd x,
\]
则成立闭式恒等
\[
\boxed{\;
\mathcal S=\sum_{j=1}^{\kappa}m_j\,\frac{\delta_j}{1+\delta_j}=S_{\mathrm{def}}.\;}
\]
并且 integrand 可改写为
\(
\frac{2\delta_j}{1+\delta_j}\partial_y U_j(x)=\frac{4\delta_j}{(x-\gamma_j)^2+(1+\delta_j)^2}
\)，与命题 \ref{prop:xi-four-rigidity-residue-depth-normalization} 锁定的 \(c=4\) 完全一致。
\end{proposition}

\begin{proof}
由
\(
U_j(x+\mathrm{i}y)
=\frac12\log\!\bigl((x-\gamma_j)^2+(y+1+\delta_j)^2\bigr)
-\frac12\log\!\bigl((x-\gamma_j)^2+(y-1-\delta_j)^2\bigr)
\)
直接求导并取 \(y=0\) 得到 \(\partial_y U_j(x)\) 的显式表达。
再利用恒等式
\[
\int_{\RR}\frac{2a}{(x-\gamma)^2+a^2}\,\dd x=2\pi\qquad(a>0),
\]
对每个 \(j\) 逐项积分即可得到
\(
\mathcal S=\sum_j m_j\frac{2\delta_j}{1+\delta_j}\cdot\frac{1}{4\pi}\cdot 2\pi
=\sum_j m_j\frac{\delta_j}{1+\delta_j}
\)，
最后一步与 \eqref{eq:xi-defect-entropy-closed-form}（命题 \ref{prop:xi-defect-entropy-closed-form-invariant}）一致。证毕。
\end{proof}

\begin{proposition}[信息可分辨性的单调收缩：Poisson 粗粒化对任意 \texorpdfstring{$f$}{f}-散度严格收缩]\label{prop:xi-poisson-coarsegraining-f-divergence-contraction}
\leanverified{paper\_xi\_poisson\_coarsegraining\_f\_divergence\_contraction}
令两组缺陷配置诱导两条非负剖面 \(H,\widetilde H\in L^1(\RR)\)，并归一化为概率密度
\(
h:=H/\|H\|_{L^1},\ \tilde h:=\widetilde H/\|\widetilde H\|_{L^1}.
\)
对任意 \(t>0\)，令 Poisson 核 \(P_t(x):=\frac{1}{\pi}\frac{t}{x^2+t^2}\)，并定义粗粒化 \(h_t:=P_t*h\)、\(\tilde h_t:=P_t*\tilde h\)。
对任意凸函数 \(f\)（且 \(f(1)=0\)）对应的 \(f\)-散度
\[
D_f(h\mid \tilde h):=\int_{\RR}\tilde h(x)\,f\!\left(\frac{h(x)}{\tilde h(x)}\right)\dd x
\]
满足数据处理单调性
\[
\boxed{\;
D_f(h_t\mid \tilde h_t)\le D_f(h\mid \tilde h)\qquad(\forall t>0).\;}
\]
若 \(f\) 严格凸，则除 \(h=\tilde h\)（制度可见域内不可区分）外，上式为严格不等式。
\end{proposition}

\begin{proof}
把卷积写成 Markov 核 \(K_t(y\mid x):=P_t(y-x)\)，则
\(
h_t(y)=\int K_t(y\mid x)h(x)\dd x
\)、\(
\tilde h_t(y)=\int K_t(y\mid x)\tilde h(x)\dd x
\)。
对每个 \(y\) 定义概率核
\[
\pi_y(\dd x):=\frac{K_t(y\mid x)\tilde h(x)}{\tilde h_t(y)}\,\dd x,
\]
则
\(
\frac{h_t(y)}{\tilde h_t(y)}=\int \frac{h(x)}{\tilde h(x)}\,\pi_y(\dd x)
\)
为 \(\frac{h}{\tilde h}\) 在 \(\pi_y\) 下的条件期望。
由 Jensen 不等式与 \(f\) 的凸性，
\[
f\!\left(\frac{h_t(y)}{\tilde h_t(y)}\right)
\le
\int f\!\left(\frac{h(x)}{\tilde h(x)}\right)\pi_y(\dd x).
\]
两边乘以 \(\tilde h_t(y)\) 并对 \(y\) 积分，交换积分次序即得
\(
D_f(h_t\mid \tilde h_t)\le D_f(h\mid \tilde h)
\)。
若 \(f\) 严格凸，则 Jensen 取等当且仅当 \(\frac{h}{\tilde h}\) 在 \(\pi_y\) 上几乎处处为常数；由于 \(K_t>0\) 在 \(\RR\times\RR\) 上处处正，这迫使 \(h=\tilde h\)。
证毕。
\end{proof}

\begin{proposition}[黑洞熵类比的可判定版本：唯一规一化使“深缺陷饱和极限”严格等于拓扑计数]\label{prop:xi-fourpi-entropy-normalization-uniqueness}
\leanverified{paper\_xi\_fourpi\_entropy\_normalization\_uniqueness}
设对任意非负剖面 \(H\in L^1(\RR)\) 定义“面积型”量 \(A(H):=\int_{\RR}H(x)\dd x\) 与“熵型”规一化 \(S_\lambda:=A(H)/\lambda\)（\(\lambda>0\) 常数）。
若要求对任意单原子缺陷 \((\gamma,\delta,m)\)：
\begin{enumerate}
  \item \(S_\lambda\) 对 Cayley 标尺重标定不变（参见注 \ref{rem:dephys-cayley-scale-gauge-four}，即只允许落在对数差/比值不变量中）；
  \item \(S_\lambda\) 对互不耦合缺陷可加；
  \item 深缺陷极限 \(\delta\to\infty\) 时 \(S_\lambda\) 严格饱和到拓扑计数 \(m\)；
\end{enumerate}
则在本文制度约定下必有且仅有
\[
\boxed{\;
\lambda=4\pi,\qquad
H(x)=\sum_{j}m_j\,\frac{4\delta_j}{(x-\gamma_j)^2+(1+\delta_j)^2},\qquad
S_\lambda(H)=\sum_{j}m_j\,\frac{\delta_j}{1+\delta_j}=S_{\mathrm{def}}.\;}
\]
\end{proposition}

\begin{proof}
由命题 \ref{prop:xi-four-rigidity-residue-depth-normalization}，残数一致化把剖面规一化唯一锁定为 \(c=4\)；从而单原子面积为
\(
A(H)=\pi\cdot 4m\,\frac{\delta}{1+\delta}
\)。
深缺陷极限要求 \(\lim_{\delta\to\infty}S_\lambda(H)=m\)，故必有 \(\lambda=4\pi\)。
另一方面，乘法伸缩半群在 bounded 坐标中的 Möbius 线性化把深度权重唯一锁定为
\(
p(\delta)=\delta/(1+\delta)
\)
（命题 \ref{prop:xi-depth-weight-uniqueness-mobius-linearization}），与 \(S_{\mathrm{def}}\) 的闭式表达 \eqref{eq:xi-defect-entropy-closed-form} 一致。
因此三条要求共同强制给出唯一规一化。证毕。
\end{proof}

\begin{remark}[系数 \texorpdfstring{$4$}{4} 的终极解释：曲率标尺、残数标尺、饱和标尺的三重一致化]\label{rem:xi-four-triple-consistency-curvature-residue-saturation}
在本文闭合制度中，常数 \(4\) 同时满足三条彼此独立的一致化条件：
\begin{enumerate}
  \item \textbf{几何侧（曲率标尺）}：圆盘双曲度量 \(ds^2_{\mathbb D}=\frac{4|dw|^2}{(1-|w|^2)^2}\) 的规一化把高斯曲率固定为 \(-1\)（命题 \ref{prop:dephys-hyperbolic-curvature-normalization-four}）；
  \item \textbf{解析侧（残数标尺）}：使 \(M_{H_c}\) 与 \(M_P\) 的残数比精确恢复无量纲深度权重（命题 \ref{prop:xi-four-rigidity-residue-depth-normalization}）；
  \item \textbf{计数侧（饱和标尺）}：使面积规一化熵在深缺陷极限精确饱和到拓扑计数（命题 \ref{prop:xi-fourpi-entropy-normalization-uniqueness}）。
\end{enumerate}
因此，“除以 \(4\)”的结构不是算术偶然，而是把边界几何量在制度允许的可观测组合中规范化为无量纲计数量的唯一机制；其与黑洞熵公式的形式相似性，严格地来源于同一类边界规一化与饱和对齐的数学结构，而非引入新的动力学解释。
\end{remark}
